% GENERATED FILE. Edit canonical lessons, then run npm run generate:books.
% canonical-source-hash: fnv1a64:5db2becbcb9f4212
% canonical-lessons: MR-C323-nirop, MR-C323-nivrtti, MR-C323-grhapraves, MR-C323-vardhapandin, MR-C323-samarop

\chapter{Farewells and Anniversaries}
\label{ch:mr-farewells-and-anniversaries}
\hlchaptermodality{323}

\begin{chapteropening}
\textbf{By the end of this chapter:} \emph{I can talk about farewells, retirement and anniversaries.}

Complete the last lesson of chapter 323: I can talk about farewells, retirement and anniversaries.
\end{chapteropening}

\section[\texorpdfstring{\mr{निरोप} (nirop)}{निरोप (nirop)}]{\mr{निरोप} (nirop) — a farewell; a message}

\label{lesson:MR-C323-nirop}



\begin{quote}
Before the new one: say the Marathi for to congratulate, then the Marathi for to inaugurate.
\end{quote}

\subsection*{You'll want to know: \mr{निरोप}}
\textbf{\mr{निरोप}} — \emph{nirop} — \textquotedblleft{}a farewell; a message\textquotedblright{}.

\begin{quote}
\mr{शाळेत} \mr{गुरुजींना} \mr{निवृत्तीच्या} \mr{दिवशी} \mr{निरोप} \mr{दिला}.
\end{quote}

\begin{quote}
\emph{śāḷet gurujīnnā nivṛttīcyā divśī nirop dilā.}
\end{quote}

\begin{quote}
At school the teacher was given a farewell on the day he retired.
\end{quote}

Masculine.

\mr{निरोप} \mr{देणे}, to see someone off; \mr{निरोप} \mr{घेणे}, to take one's leave.

\begin{cousinweb}
\textbf{\mr{निरोप}} (\emph{nirop}) is the Marathi word for a message and a leave-taking; its origin is not settled.
\end{cousinweb}

\begin{grammarlens}[title={What you've built}]
One more word for telling what happened.
\end{grammarlens}

\subsection*{Your turn}
\begin{itemize}
\raggedright
  \item \emph{Say it:} \emph{nirop}
  \item \emph{Say it:} \emph{nirop}, once more
  \item \emph{Say it:} \emph{śāḷet gurujīnnā nivṛttīcyā divśī nirop dilā.} — At school the teacher was given a farewell on the day he retired.
  \item \emph{Recall:} say the Marathi for to congratulate, then the Marathi for to inaugurate, then say \emph{nirop} again
\end{itemize}

\begin{tcolorbox}[breakable,colback=teal!4,colframe=teal!35!black,title={Before you move on}]
What does \mr{निरोप} mean? (\textquotedblleft{}A farewell; a message\textquotedblright{}.) Say it once more.
\end{tcolorbox}
\section[\texorpdfstring{\mr{निवृत्ती} (nivṛttī)}{निवृत्ती (nivṛttī)}]{\mr{निवृत्ती} (nivṛttī) — retirement}

\label{lesson:MR-C323-nivrtti}



\begin{quote}
Before the new one: say the Marathi for to inaugurate, then the Marathi for a farewell.
\end{quote}

\subsection*{You'll want to know: \mr{निवृत्ती}}
\textbf{\mr{निवृत्ती}} — \emph{nivṛttī} — \textquotedblleft{}retirement\textquotedblright{}.

\begin{quote}
\mr{निवृत्तीनंतर} \mr{बाबा} \mr{रोज} \mr{बागेत} \mr{काम} \mr{करतात}.
\end{quote}

\begin{quote}
\emph{nivṛttīnantar bābā roj bāget kām kartāt.}
\end{quote}

\begin{quote}
Since retiring Dad works in the garden every day.
\end{quote}

Feminine.

\mr{निवृत्तीनंतर}, after retirement.

\begin{cousinweb}
\textbf{\mr{निवृत्ती}} (\emph{nivṛttī}) is a tatsama, taken from Sanskrit \emph{nivṛtti}, 'a turning back, ceasing from work' (Monier-Williams).
\end{cousinweb}

\begin{grammarlens}[title={What you've built}]
One more word for telling what happened.
\end{grammarlens}

\subsection*{Your turn}
\begin{itemize}
\raggedright
  \item \emph{Say it:} \emph{nivṛttī}
  \item \emph{Say it:} \emph{nivṛttī}, once more
  \item \emph{Say it:} \emph{nivṛttīnantar bābā roj bāget kām kartāt.} — Since retiring Dad works in the garden every day.
  \item \emph{Recall:} say the Marathi for to inaugurate, then the Marathi for a farewell, then say \emph{nivṛttī} again
\end{itemize}

\begin{tcolorbox}[breakable,colback=teal!4,colframe=teal!35!black,title={Before you move on}]
What does \mr{निवृत्ती} mean? (\textquotedblleft{}Retirement\textquotedblright{}.) Say it once more.
\end{tcolorbox}
\section[\texorpdfstring{\mr{गृहप्रवेश} (gṛhapraveś)}{गृहप्रवेश (gṛhapraveś)}]{\mr{गृहप्रवेश} (gṛhapraveś) — a housewarming}

\label{lesson:MR-C323-grhapraves}



\begin{quote}
Before the new one: say the Marathi for a farewell, then the Marathi for retirement.
\end{quote}

\subsection*{You'll want to know: \mr{गृहप्रवेश}}
\textbf{\mr{गृहप्रवेश}} — \emph{gṛhapraveś} — \textquotedblleft{}a housewarming\textquotedblright{}.

\begin{quote}
\mr{नव्या} \mr{घराचा} \mr{गृहप्रवेश} \mr{रविवारी} \mr{झाला}.
\end{quote}

\begin{quote}
\emph{navyā gharācā gṛhapraveś ravivārī jhālā.}
\end{quote}

\begin{quote}
The housewarming of the new house was on Sunday.
\end{quote}

Masculine.

\mr{गृह}, a house, + \mr{प्रवेश}, an entry.

\begin{cousinweb}
\textbf{\mr{गृहप्रवेश}} (\emph{gṛhapraveś}) is a tatsama, taken from Sanskrit \emph{gṛhapraveśa}, 'a solemn entrance into a house' (Monier-Williams).
\end{cousinweb}

\begin{grammarlens}[title={What you've built}]
One more word for telling what happened.
\end{grammarlens}

\subsection*{Your turn}
\begin{itemize}
\raggedright
  \item \emph{Say it:} \emph{gṛhapraveś}
  \item \emph{Say it:} \emph{gṛhapraveś}, once more
  \item \emph{Say it:} \emph{navyā gharācā gṛhapraveś ravivārī jhālā.} — The housewarming of the new house was on Sunday.
  \item \emph{Recall:} say the Marathi for a farewell, then the Marathi for retirement, then say \emph{gṛhapraveś} again
\end{itemize}

\begin{tcolorbox}[breakable,colback=teal!4,colframe=teal!35!black,title={Before you move on}]
What does \mr{गृहप्रवेश} mean? (\textquotedblleft{}A housewarming\textquotedblright{}.) Say it once more.
\end{tcolorbox}
\section[\texorpdfstring{\mr{वर्धापनदिन} (vardhāpandin)}{वर्धापनदिन (vardhāpandin)}]{\mr{वर्धापनदिन} (vardhāpandin) — an anniversary}

\label{lesson:MR-C323-vardhapandin}



\begin{quote}
Before the new one: say the Marathi for retirement, then the Marathi for a housewarming.
\end{quote}

\subsection*{You'll want to know: \mr{वर्धापनदिन}}
\textbf{\mr{वर्धापनदिन}} — \emph{vardhāpandin} — \textquotedblleft{}an anniversary\textquotedblright{}.

\begin{quote}
\mr{आज} \mr{आमच्या} \mr{कंपनीचा} \mr{पन्नासावा} \mr{वर्धापनदिन} \mr{आहे}.
\end{quote}

\begin{quote}
\emph{āj āmacyā kampnīcā pannāsāvā vardhāpandin āhe.}
\end{quote}

\begin{quote}
Today is our company's fiftieth anniversary.
\end{quote}

Masculine.

\mr{वर्धापन}, a birthday rite, + \mr{दिन}, a day.

\begin{cousinweb}
\textbf{\mr{वर्धापनदिन}} (\emph{vardhāpandin}) is built from Sanskrit \emph{vardhāpana}, a rite kept on a birthday (Monier-Williams), and \emph{dina}, 'a day'.
\end{cousinweb}

\begin{grammarlens}[title={What you've built}]
One more word for telling what happened.
\end{grammarlens}

\subsection*{Your turn}
\begin{itemize}
\raggedright
  \item \emph{Say it:} \emph{vardhāpandin}
  \item \emph{Say it:} \emph{vardhāpandin}, once more
  \item \emph{Say it:} \emph{āj āmacyā kampnīcā pannāsāvā vardhāpandin āhe.} — Today is our company's fiftieth anniversary.
  \item \emph{Recall:} say the Marathi for retirement, then the Marathi for a housewarming, then say \emph{vardhāpandin} again
\end{itemize}

\begin{tcolorbox}[breakable,colback=teal!4,colframe=teal!35!black,title={Before you move on}]
What does \mr{वर्धापनदिन} mean? (\textquotedblleft{}An anniversary\textquotedblright{}.) Say it once more.
\end{tcolorbox}
\section[\texorpdfstring{\mr{समारोप} (samārop)}{समारोप (samārop)}]{\mr{समारोप} (samārop) — a closing; a conclusion}

\label{lesson:MR-C323-samarop}



\begin{quote}
Before the new one: say the Marathi for a housewarming, then the Marathi for an anniversary.
\end{quote}

\subsection*{You'll want to know: \mr{समारोप}}
\textbf{\mr{समारोप}} — \emph{samārop} — \textquotedblleft{}a closing; a conclusion\textquotedblright{}.

\begin{quote}
\mr{संमेलनाचा} \mr{समारोप} \mr{एका} \mr{गाण्याने} \mr{झाला}.
\end{quote}

\begin{quote}
\emph{sammelanācā samārop ekā gāṇyāne jhālā.}
\end{quote}

\begin{quote}
The gathering closed with a song.
\end{quote}

Masculine.

\mr{समारोप} \mr{करणे}, to close an event.

\begin{cousinweb}
\textbf{\mr{समारोप}} (\emph{samārop}) is a tatsama, taken from Sanskrit \emph{samāropa}, 'a placing upon' (Monier-Williams); the Marathi sense, the close of an event, is its own.
\end{cousinweb}

\begin{grammarlens}[title={What you've built}]
One more word for telling what happened.
\end{grammarlens}

\subsection*{Your turn}
\begin{itemize}
\raggedright
  \item \emph{Say it:} \emph{samārop}
  \item \emph{Say it:} \emph{samārop}, once more
  \item \emph{Say it:} \emph{sammelanācā samārop ekā gāṇyāne jhālā.} — The gathering closed with a song.
  \item \emph{Recall:} say the Marathi for a housewarming, then the Marathi for an anniversary, then say \emph{samārop} again
\end{itemize}

\begin{tcolorbox}[breakable,colback=teal!4,colframe=teal!35!black,title={Before you move on}]
What does \mr{समारोप} mean? (\textquotedblleft{}A closing; a conclusion\textquotedblright{}.) Say it once more.
\end{tcolorbox}
